\documentclass[11pt]{article}
\usepackage[margin=1in]{geometry}
\usepackage{amsmath,amssymb,amsthm}
\usepackage[colorlinks=true,linkcolor=blue,citecolor=blue,urlcolor=blue]{hyperref}

\newtheorem{theorem}{Theorem}
\newtheorem{remark}[theorem]{Remark}

\title{A Counterexample to Vasilyev's Cycle Ratio Inequality $2^\ell < 3^{a+1}$ in the Up-Down Graph}
\author{Salem Levinstein}
\date{November 26, 2026}

\begin{document}
\maketitle

\begin{abstract}
For any subset $S \subseteq \mathbb{N} = \{1,2,\dots\}$, Vasilyev (\href{https://arxiv.org/abs/2611.20344}{arXiv:2611.20344v2}) studies the Collatz-type map $T_S(x) = \lceil 3x/2 \rceil$ for $x \in S$ and $T_S(x) = \lceil x/2 \rceil$ for $x \notin S$, whose primitive cycles correspond to simple directed cycles of the up-down graph $G$. In Question~7.1, Vasilyev asks whether every cycle of length $\ell$ with $a$ up-steps satisfies $2^\ell < 3^{a+1}$ (equivalently $a = \lfloor \ell \log_3 2 \rfloor$), having verified this inequality for all cycles with $\ell \le 30$ and all $1{,}133{,}265{,}572$ cycles contained in $[1,160]$. We answer this question in the negative by exhibiting an explicit simple cycle of length $\ell = 382$ with $a = 240$ up-steps, for which $3^{241} < 2^{382}$ and $2^{382}/3^{240} \approx 3.050403 > 3$.
\end{abstract}

\section{Introduction}
Let $\mathbb{N} = \{1,2,3,\dots\}$, and define the upward and downward step maps
\[
U(x) = \left\lceil \frac{3x}{2} \right\rceil = \left\lfloor \frac{3x+1}{2} \right\rfloor,
\qquad
D(x) = \left\lceil \frac{x}{2} \right\rceil = \left\lfloor \frac{x+1}{2} \right\rfloor.
\]
For any subset $S \subseteq \mathbb{N}$, define $T_S : \mathbb{N} \to \mathbb{N}$ by $T_S(x) = U(x)$ if $x \in S$ and $T_S(x) = D(x)$ if $x \notin S$. Let $G$ be the directed \emph{up-down graph} on $\mathbb{N}$ with edges $x \to U(x)$ and $x \to D(x)$. By the Dictionary proposition of \cite{Vasilyev2026}, primitive cycles of maps $T_S$ are in bijection with simple directed cycles of $G$.

For any cycle $c = (x_0,\dots,x_{\ell-1})$ of length $\ell$ with $a$ up-steps and $\ell - a$ down-steps, Theorem~4.1(c) of \cite{Vasilyev2026} establishes $3^a < 2^\ell$ via the exact product identity
\begin{equation}\label{eq:prod}
\frac{2^\ell}{3^a} = \prod_{i=0}^{\ell-1} \left(1 + \frac{e_i}{c_i x_i}\right),
\end{equation}
where $e_i = x_i \bmod 2 \in \{0,1\}$, $c_i = 3$ if $x_{i+1}=U(x_i)$, and $c_i = 1$ if $x_{i+1}=D(x_i)$. In Section~7, Vasilyev poses \textbf{Question~7.1}: does every cycle satisfy $2^\ell < 3^{a+1}$ (i.e., $a = \lfloor \ell \log_3 2 \rfloor$)? Vasilyev proved this for single-run cycles (Theorem~4.1(d)) and verified it for all $\ell \le 30$ and all $1{,}133{,}265{,}572$ cycles whose elements lie in $[1,160]$, among which the maximum ratio $2^\ell/3^a$ is $2.398$ (at $\ell = 71$).

\section{Explicit Counterexample}

\begin{theorem}\label{thm:main}
Question~7.1 of \cite{Vasilyev2026} has a negative answer: there exists a subset $S \subset \mathbb{N}$ and a primitive cycle $c = (x_0,\dots,x_{381})$ of $T_S$ with length $\ell = 382$ and $a = 240$ up-steps satisfying
\[
3^{241} < 2^{382}.
\]
In particular, $2^{382}/3^{240} \approx 3.050403 > 3$ and $a = 240 = \lfloor 382 \log_3 2 \rfloor - 1 \neq \lfloor \ell \log_3 2 \rfloor$.
\end{theorem}

\begin{proof}
Consider the sequence $c = (x_0,x_1,\dots,x_{381})$ of $382$ integers $\ge 2$ listed in full in Table~\ref{tab:cycle}. Its vertices are pairwise distinct, ranging from $2$ to $258{,}477{,}331{,}094$, and partition into $13$ alternating up-runs and down-runs whose local minima (up-run starting vertices) are
\[
2,\; 6,\; 10,\; 16,\; 20,\; 29,\; 33,\; 38,\; 40,\; 42,\; 46,\; 51,\; 56.
\]
For every index $i \in \{0,\dots,381\}$ (modulo $382$), direct arithmetic verifies $x_{i+1} \in \{U(x_i), D(x_i)\}$. Since $x_i \ge 2$ implies $D(x_i) < x_i < U(x_i)$, the step type at every vertex is uniquely determined: $a = 240$ indices satisfy $x_{i+1} = U(x_i)$ and $142$ satisfy $x_{i+1} = D(x_i)$. Setting $S = \{x_i : x_{i+1} = U(x_i)\}$, pairwise distinctness implies $T_S(x_i) = x_{i+1}$ for all $i$, so $c$ is a primitive cycle of $T_S$ of exact period $382$.

Finally, exact integer exponentiation gives
\begin{align*}
3^{241} &= 968772506150712031215849558055091779642661983952325151332908\\
        &\qquad\quad 6812944808568617534785161369767311066873140028406651043,\\
2^{382} &= 985050154909861980306976002503590345126993481761636166698707\\
        &\qquad\quad 3351061430442874302652853566563721228910201656997576704,
\end{align*}
whence $3^{241} < 2^{382} < 3^{242}$, proving $2^\ell > 3^{a+1}$ and $a = \lfloor \ell \log_3 2 \rfloor - 1$.
\end{proof}

\begin{table}[h!]
\centering\footnotesize
\begin{minipage}{0.96\textwidth}
\raggedright\ttfamily
2, 3, 5, 8, 12, 18, 27, 41, 62, 93, 140, 210, 315, 473, 710, 1065, 1598, 2397, 3596, 5394, 2697, 1349, 675, 338, 169, 85, 43, 22, 11, 6, 9, 14, 21, 32, 48, 72, 108, 162, 243, 365, 548, 822, 1233, 1850, 2775, 4163, 6245, 9368, 4684, 2342, 1171, 586, 293, 147, 74, 37, 19, 10, 15, 23, 35, 53, 80, 120, 180, 270, 405, 608, 912, 1368, 2052, 3078, 4617, 6926, 10389, 15584, 23376, 35064, 52596, 78894, 118341, 177512, 266268, 399402, 599103, 898655, 1347983, 2021975, 3032963, 4549445, 6824168, 10236252, 15354378, 23031567, 34547351, 51821027, 77731541, 116597312, 174895968, 262343952, 393515928, 590273892, 885410838, 1328116257, 1992174386, 2988261579, 4482392369, 6723588554, 10085382831, 15128074247, 22692111371, 34038167057, 51057250586, 76585875879, 114878813819, 172318220729, 258477331094, 129238665547, 64619332774, 32309666387, 16154833194, 8077416597, 4038708299, 2019354150, 1009677075, 504838538, 252419269, 126209635, 63104818, 31552409, 15776205, 7888103, 3944052, 1972026, 986013, 493007, 246504, 123252, 61626, 30813, 15407, 7704, 3852, 1926, 963, 482, 241, 121, 61, 31, 16, 24, 36, 54, 81, 122, 183, 275, 413, 620, 310, 155, 78, 39, 20, 30, 45, 68, 102, 153, 230, 115, 58, 29, 44, 66, 99, 149, 224, 336, 504, 756, 1134, 1701, 2552, 3828, 5742, 8613, 12920, 19380, 29070, 43605, 65408, 98112, 147168, 220752, 331128, 496692, 745038, 1117557, 1676336, 2514504, 3771756, 5657634, 8486451, 4243226, 2121613, 1060807, 530404, 265202, 132601, 66301, 33151, 16576, 8288, 4144, 2072, 1036, 518, 259, 130, 65, 33, 50, 75, 38, 57, 86, 129, 194, 291, 437, 656, 984, 1476, 2214, 3321, 4982, 7473, 11210, 16815, 25223, 37835, 56753, 85130, 127695, 191543, 287315, 430973, 646460, 323230, 161615, 80808, 40404, 20202, 10101, 5051, 2526, 1263, 632, 316, 158, 79, 40, 60, 90, 135, 203, 305, 458, 687, 1031, 1547, 2321, 3482, 5223, 7835, 11753, 17630, 26445, 39668, 59502, 89253, 133880, 200820, 301230, 451845, 677768, 338884, 169442, 84721, 42361, 21181, 10591, 5296, 2648, 1324, 662, 331, 166, 83, 42, 63, 95, 143, 215, 323, 485, 728, 364, 182, 91, 46, 69, 104, 156, 234, 351, 527, 791, 1187, 1781, 2672, 4008, 6012, 9018, 13527, 20291, 30437, 45656, 68484, 102726, 51363, 25682, 12841, 6421, 3211, 1606, 803, 402, 201, 101, 51, 77, 116, 174, 261, 392, 588, 882, 441, 221, 111, 56, 84, 126, 189, 284, 426, 639, 959, 1439, 2159, 3239, 4859, 7289, 10934, 16401, 24602, 12301, 6151, 3076, 1538, 769, 385, 193, 97, 49, 25, 13, 7, 4
\end{minipage}
\caption{The 382-vertex simple cycle in the up-down graph $G$ with $a = 240$ up-steps and $2^{382}/3^{240} \approx 3.050403 > 3$.}\label{tab:cycle}
\end{table}

\begin{remark}
The follow-up part of Question~7.1 asks for $\sup 2^\ell / 3^a$ over all simple cycles of $G$. Theorem~\ref{thm:main} establishes $\sup 2^\ell / 3^a \ge 2^{382}/3^{240} \approx 3.0504$, and extending the construction to more small odd targets yields larger ratios (e.g., $\approx 3.72$ at $\ell = 12{,}913$). Determining whether $\sup 2^\ell/3^a = \infty$ remains open.
\end{remark}

\section*{Formal Verification}
Theorem~\ref{thm:main} is formally verified in Lean~4 (\texttt{Conjectures/VasilyevUpDownCycles.lean}, theorem \texttt{vasilyev\_question\_7\_1\_disproof}) with zero \texttt{sorry} and standard axioms only, and independently verified in exact integer and rational arithmetic by \texttt{verify.py}.

\begin{thebibliography}{9}
\bibitem{Vasilyev2026}
O.~Vasilyev, \emph{Up or down at every integer: Collatz-type maps of arbitrary sets and the up-down graph}, arXiv:2611.20344v2 [math.NT], November 2026.
\end{thebibliography}
\end{document}
